\documentclass[11pt]{article}
\usepackage{fontspec}
\usepackage[french]{babel}
\usepackage{amsmath}
\usepackage{amssymb}
\usepackage{amsthm}
\usepackage{mathrsfs}
\usepackage{tikz-cd}
\usepackage{enumitem}
\usetikzlibrary{babel}

\newcommand{\printedpage}[1]{\marginpar{\footnotesize p.~#1}}
\newcommand{\msp}[1]{\marginpar{\footnotesize [#1]}}

\newcommand{\Hom}{\mathrm{Hom}}
\newcommand{\Fl}{\mathrm{Fl}}
\newcommand{\Ob}{\mathrm{Ob}}
\newcommand{\Cat}{\mathrm{Cat}}
\newcommand{\Ens}{\mathrm{Ens}}
\newcommand{\pr}{\mathrm{pr}}
\newcommand{\Ker}{\mathrm{Ker}}
\newcommand{\Coker}{\mathrm{Coker}}
\newcommand{\mon}{\mathrm{mon}}
\newcommand{\deq}{\mathrel{\overset{\text{déf}}{=}}}

\begin{document}

\printedpage{}
% no page number printed (title page)
% head: (none)

\begin{center}
{\LARGE\textsf{LES DÉRIVATEURS}}

\bigskip

\textbf{ALEXANDRE GROTHENDIECK}

\bigskip\bigskip

{\large\textbf{Appendice au Chapitre XII}}
\end{center}

\newpage

\printedpage{1}
% head: A. Grothendieck : Dérivateurs, Ch. XII, Édition des Universités de Montpellier II et Paris VII
% foot: Version 20 juin 2003

\msp{page 1}
\textbf{Proposition} (?) (\footnote{Cf. pages 11 ff. pour les contrexemples-clef.}). \textit{Soit $X$ dans} $\Cat$\textit{, avec $b$, $c$ dans $X$ tels que}

\begin{itemize}
\item[\textit{a)}] \textit{$b\times c$ est dans $X$.}
\item[\textit{b)}] \textit{$\forall$ $x$ dans $X$, $b(x)$ ($\deq$} $\Hom_X(x,b)$\textit{) et $c(x)$ sont $\neq\emptyset$, i.e. $X/b \longrightarrow X$ et $X/c \longrightarrow X$ sont surjectifs sur les objets.}
\end{itemize}

\textit{Alors $X$ est $\mathbf{W}_\infty$-asphérique.} (\footnote{Faux, cf. plus bas. $X$ est 1-connexe, mais pas nécessairement 2-connexe.})

Le fait que $b$ soit un plus grand élément de $\Ob X$ assure déjà la 0-connexité.

Pour la 1-connexité, prouvons que dans le groupoïde fondamental $\Pi(X)$, $a \deq b\times c$ devient un élément final.
\[
\begin{tikzcd}
b\times c \arrow[r] \arrow[d] & b \\
c &
\end{tikzcd}
\]

Pour ceci, prouvons d'abord que $\forall$ $x$ dans $X$, tous les éléments de $a(x) = b(x)\times c(x)$ sont égaux dans $\Fl\Pi(X)$. En effet, comme $b\times c \longrightarrow b$ est \texttt{[un]} isomorphisme, donc \texttt{[un]} monomorphisme dans $\Pi(X)$, cela montre déjà que si $\beta,\beta'\in b(x)$, $\gamma,\gamma'\in c(x)$, d'où $(\beta,\gamma)$, $(\beta,\gamma')$, $(\beta',\gamma')\in b(x)\times c(x)$, on a $(\beta,\gamma) = (\beta,\gamma')$ dans $\Pi(X)$ (puisque $\pr_1(\beta,\gamma) = \pr_1(\beta,\gamma')$, et de même (en utilisant que $\pr_2$ est \texttt{[un]} isomorphisme, donc \texttt{[un]} monomorphisme) $(\beta,\gamma') = (\beta',\gamma')$). Donc tous les $(\beta,\gamma)\in b(x)\times c(x) = a(x)$ sont égaux dans $\Fl\Pi(X)$. Pour achever de prouver que $a$ est objet final de $\Pi(X)$, il suffit de voir que toute flèche \textit{dans} $\Pi(X)$ d'un $x\in\Ob X$ dans $a$, provient d'une flèche $x\longrightarrow a$ dans $X$. Mais une flèche dans $\Pi(X)$ est obtenue par un chemin dans $X$
\[
\begin{tikzcd}
x = x_0 \arrow[r, "u_1"] \arrow[dr, "\alpha_0"] & x_1 \arrow[d, "\alpha_1"'] & x_2 \arrow[l, "u_2"'] \arrow[r] \arrow[dl, "\alpha_2"'] & {} & \cdots & {} & x_{2N} = a \arrow[l, "u_{2N}"'] \\
& a. & & & & &
\end{tikzcd}
\]

Prenons une flèche $\alpha_0 : x\longrightarrow a_0$ dans $X$, je dis que la % [sic] a_0 -> a, the arrows α_i go to a (no a_0 is defined)
\msp{page 2}
classe d'homotopie du chemin $c$ (i.e. son image dans $\Fl(\Pi(X))$ est égale à celle de $u_0$. % [sic] u_0 -> \alpha_0 (?), no u_0 is defined; also the parenthesis opened at "(i.e." is not closed
Pour le voir, choisissons pour tout $i\in[0,2N]$ une flèche $\alpha_i : x_i\longrightarrow a$. (Ça existe grâce à l'hypothèse b)). Je dis que pour chaque $i\in[0,2N]$, le triangle
\printedpage{2}
% head: A. Grothendieck : Dérivateurs, Ch. XII, Édition des Universités de Montpellier II et Paris VII
% foot: Version 20 juin 2003
\[
\begin{tikzcd}[column sep=small]
x_0 \arrow[rr, "c_i"] \arrow[dr, "\alpha_0"] & & x_i \arrow[dl, "\alpha_i"'] \\
& a &
\end{tikzcd}
\]
dans $\Pi(X)$ est commutatif, où $c_i$ désigne la classe d'homotopie du chemin de $x_0$ à $x_i$ définie par les $u_i$ ($i\in[1,i]$) \texttt{[plutôt} ($i\in[1,2N]$)\texttt{]}. Nous le faisons par récurrence sur $i$, l'assertion est triviale si $i = 0$ ($c_0$ représente l'application identitique). Le diagramme % [sic] identitique -> identique, non-existent word
\[
\begin{tikzcd}[column sep=small]
x_0 \arrow[rr, "c_i"] \arrow[dr, "\alpha_0"] \arrow[rrrr, bend left=15, "c_{i+1}"] & & x_i \arrow[dl, "\alpha_i"'] \arrow[rr, "u_{i+1}"] & & x_{i+1} \arrow[dlll, "\alpha_{i+1}"'] \\ % ?? uncertain: the curved arrow c_{i+1} -- solid or dashed
& a & & &
\end{tikzcd}
\]
montre qu'il suffit de prouver que pour $i\in[0,2N-1]$ le triangle
\[
\begin{tikzcd}[column sep=small]
x_i \arrow[rr, "u_{i+1}"] \arrow[dr, "\alpha_i"'] & & x_{i+1} \arrow[dl, "\alpha_{i+1}"] \\
& a &
\end{tikzcd}
\]
est commutatif dans $\Pi(X)$. (\textbf{NB} $u_{i+1}$ est en sens directe ou rétrograde, suivant que % [sic] directe -> direct, agreement with "sens"
$i$ est paire ou impaire.) Mais cela résulte du fait que l'image de $\Hom_X(x_i,a)$ (resp. $\Hom_X(x_{i+1},a)$) % [sic] paire ou impaire -> pair ou impair, agreement with "i"
\msp{page 3}
dans $\Fl\Pi(X)$ est réduite à un seul élément.

Ainsi on trouve que $X$ \textit{est} 1-\textit{connexe.} Il y a donc de l'espoir! Il reste à prouver
\[
H^i(X,M) = 0, \qquad i\geq 2, \quad M \text{ coefficient constant.}
\]

J'ai regardé le cas lorsqu'on a ceci : $b = c$, $\Ob X = \{b, b\times b\}$, $\Hom_X(b\times b)$ \texttt{[plutôt} $\Hom(b\times b, b)$\texttt{]} $= \{$ réunion disjointe des ensembles $u\circ\pr_1$ et $u\circ\pr_2$, avec $u\in\Hom_X(b,b) = G$ un monoïde quelconque $\}$. Alors $X^\circ$ est fibrée, via $I\longmapsto b^I$, sur la sous-catégorie pleine $S$ de $\Ens$ formée de $\{0\}$ et $\{0,1\}$, et donc $X$ est cofibrée (en particulier propre) sur $S^\circ \deq T$. (NB $S^\circ = T$ a un objet initial, mais non un objet final. Sa cohomologie à valeurs dans un préfaisceau quelconque n'est pas toujours triviale, cf. ci-dessous.)

On trouve \texttt{[`On trouve' barré]}
\[
H^n(X,M) \Longleftarrow H^p(T, \underbrace{R^q f_*(M_X)}_{= \mathcal{H}^q}), \qquad \text{où } f : X \longrightarrow S^\circ
\]
\texttt{[}$n = p+q$\texttt{]}. J'admet ici (provisoirement) ceci : % [sic] J'admet -> J'admets, verb ending
\[
\begin{aligned}
H^p(T,F) &= 0 \ , &&\text{pour } p\geq 2, \text{ tout } F\in T^\wedge \ ,\\
H^0(T,F) &= \Ker\delta \ , &&\\
H^2(T,F) &= \Coker\delta \ , && % [sic] H^2(T,F) -> H^1(T,F), the first line gives H^p(T,F) = 0 for p >= 2
\end{aligned}
\]
\printedpage{3}
% head: A. Grothendieck : Dérivateurs, Ch. XII, Édition des Universités de Montpellier II et Paris VII
% foot: Version 20 juin 2003
où $\delta : F_0 \longrightarrow F_1$ est défini comme la différence $\delta = p_1^* - p_2^*$, où $p_1$, $p_2$ sont les deux projections de $\{0,1\}\simeq\{0\}^2$ dans $T = S^\circ$. On a de plus
\[
(\mathcal{H}^q)_0 = H^q(G,M), \qquad (\mathcal{H}^q)_1 = H^q(G\times G, M),
\]
et
\[
p_1^*,\ p_2^*\ :\ H^q(G,M) \rightrightarrows H^q(G{\times}G, M) \qquad \text{sont les transposés de } \pr_1,\ \pr_2 : G\times G\longrightarrow G.
\]

\msp{page 4}
On vérifie aisément que
\[
H^0(T,\mathcal{H}^q) = 0 \qquad \text{si } q\geq 1,
\]
et il reste
\[
H^n(X,M) \simeq H^1(S,\mathcal{H}^{n-1}) = \Coker(\delta : H^{n-1}(G) \longrightarrow H^{n-1}(G\times G)). % [sic] H^1(S, ...) -> H^1(T, ...), the cohomology is taken on T = S°
\]
Cela donne, pour $n = 2$,
\[
\begin{aligned}
H^2(X,M) &\simeq \Coker\Big(\delta : \Hom_\mon(G,M) \longrightarrow \underbrace{\mathit{Hom}_\mon(G\times G, M)}_{\Hom_\mon(G,M)\times\Hom_\mon(G,M)} \ \Big)\\
&\simeq \Coker\Big(u\longmapsto (u,-u) : H \longrightarrow H\times H\Big),
\end{aligned}
\]
où $H$ est le groupe abélien $\Hom_\mon(G,M)$.

D'où
\[
H^2(X,M) \ \simeq\ H = \Hom_\mon(G,M).
\]
Donc
\begin{quote}
$X$ est 2-connexe si et seulement si le groupe \textit{abélien} enveloppant du monoïde $G$ est nul.
\end{quote}

Ainsi, si $G$ est un groupe abélien $\neq 0$, $X$ \textit{n'est pas} 2-connexe.

\textbf{Corollaire.} $\Psi = {}$\begin{tikzcd}[ampersand replacement=\&, row sep=tiny, column sep=small, baseline=(A.base)]
\& b_0 \\
|[alias=A]| a_0 \arrow[ur] \arrow[dr] \& \\
\& c_0
\end{tikzcd}
\textit{n'est} pas \textit{totalement} $\mathbf{W}_\infty$\textit{-asphérique; i.e. si on a un diagramme} \begin{tikzcd}[ampersand replacement=\&, row sep=tiny, column sep=small, baseline=(B.base)]
\& b \\
|[alias=B]| a \arrow[ur] \arrow[dr] \& \\
\& c
\end{tikzcd}
\textit{dans une catégorie} $X\in\Ob\Cat$\textit{, tel que sa somme amalgamée dans $X^\wedge$ soit égale à $e_X$, il ne s'ensuit pas que $X$ soit $\mathbf{W}_\infty$-asphérique.} (Mais peut-être voudra bien être 1-connexe?)

\printedpage{4}
% head: A. Grothendieck : Dérivateurs, Ch. XII, Édition des Universités de Montpellier II et Paris VII
% foot: Version 20 juin 2003

\msp{page 5}
\textbf{Question.} \textit{Y-a-t-il d'autres objets totalement} $\mathbf{W}_\infty$\textit{-asphériques de} $\Cat$\textit{, que les catégories ayant un objet final?} Je suspecte à présent que non -- donc que l'on ne peut améliorer substantiellement sur la notion déjà connue d'immersion ouverte de Dwyer.

Dans le diagramme
\[
\begin{tikzcd}
U \arrow[r, hook, "i"] \arrow[d, "p"'] \arrow[dr, phantom, "\substack{\text{cocart.}\\ \text{et cart.}}"] & X \arrow[d, "q"'] & T \arrow[l, hook', dotted, "j"'] \arrow[d, dotted, "\text{iso}"] \\
U' \arrow[r, hook, "i'"] & X' & T', \arrow[l, hook', dotted, "j'"']
\end{tikzcd}
\tag*{$(*)$}
\]
prenons
\[
\begin{aligned}
U = \Psi, \quad & T = \{t\} = e \text{ catégorie finale,}\\
& t \text{ objet final de } X,\\
& \text{i.e. } H_t = (u\longmapsto \Hom_X(u,t)) = e_U \in \Ob U^\wedge \ \mbox{(\footnotemark)}.
\end{aligned}
\]
\footnotetext{i.e. $X$ déduit de $U = \Psi$ en lui rajoutant un objet final $t$ -- donc $X\simeq\boldsymbol{\Delta}^1\times\boldsymbol{\Delta}^1$,
\[
\begin{tikzcd}[ampersand replacement=\&,
execute at end picture={
\draw[rounded corners] ([xshift=-10pt,yshift=10pt]A.center) -- ([xshift=-10pt,yshift=-10pt]C.center) -- ([xshift=10pt,yshift=-10pt]D.center) -- ([xshift=10pt,yshift=3pt]D.center) -- ([xshift=3pt,yshift=3pt]C.center) -- ([xshift=3pt,yshift=10pt]A.center) -- cycle;
\node (P) at ([xshift=-35pt,yshift=-22pt]C.center) {$\Psi$};
\draw[-stealth] (P) -- ([xshift=-11pt,yshift=-4pt]C.center);
}]
|[alias=A]| {} \arrow[r] \& {} \\
|[alias=C]| {} \arrow[u] \arrow[r] \& |[alias=D]| {} \arrow[u]
\end{tikzcd}
\ .
\]}
Prenons pour $U'$ la catégorie dont il est question dans la `proposition' hypothétique (et fausse) p. 1 avec $b = c$, plus précisément une catégorie avec deux objets $b$ et $b\times b$, un groupe (voire un monoïde) $G$ d'automorphismes \texttt{[(voire d'endomorphismes)]} de $b$, et les autres flèches celles qui s'en déduisent, en exigeant que $b\times b$ soit bel et bien le carré cartésien de $b$. On trouve une catégorie cofibrée (donc propre) sur la catégorie $T$ opposée de celle des ensembles finis de la forme $\{0\}$ et $\{0,1\}$ (deux objets) avec comme fibres $B_G$ et $B_{G\times G}$ respectivement ($B_\pi =$ catégorie avec un seul objet, et monoïde d'endomorphismes $\pi$). Le foncteur
\[
U = \Psi \longrightarrow U'
\]
est celui donné par le diagramme

\msp{page 6}
\[
\begin{tikzcd}[row sep=small]
& b \\
b\times b \arrow[ur, "\pr_1"] \arrow[dr, "\pr_2"'] & \\
& b
\end{tikzcd}
\]

\end{document}
